\chapter{System Overview}
\label{ch:system-overview}

The system architecture of QCanvas follows a modular layered design separating conversion, validation, simulation, and UI concerns. This design ensures scalability, maintainability, and flexibility while supporting the hybrid CPU--QPU execution model \cite{peruzzo2014variational}. The QSim SDK operates as an independently deployable execution kernel that is integrated into the QCanvas backend while also being accessible directly via PyPI.

\section{Architectural Design}

QCanvas employs a hybrid architecture where the Next.js frontend handles UI and simple operations, while the FastAPI backend manages complex quantum computations. QCanvas orchestrates compilation and validation; QSim executes circuits on simulation backends. Both the \texttt{qcanvas-sdk} and \texttt{qsim-sdk} are published on PyPI, making their capabilities accessible to the broader Python community beyond the Web IDE.

\subsection{Layered Responsibilities}
\begin{itemize}
    \item \textbf{Presentation Layer (Next.js)}: Monaco editor UI, D3.js circuit visualization, Chart.js results dashboards, gamification and community UI, and client-side state management via Zustand.
    \item \textbf{API Layer (FastAPI)}: REST endpoints for conversion/validation/simulation, WebSocket for progress streams, structured error handling, request validation, authentication, and rate limiting.
    \item \textbf{Conversion Layer (QCanvas SDK)}: AST parsing, intelligent gate mapping, OpenQASM 3.0 generation via \texttt{QASM3Builder}, validation, conversion analytics, and logging.
    \item \textbf{Execution Layer (QSim SDK)}: Simulation backends (statevector, density matrix, stabilizer), job scheduling, result aggregation, caching, and metrics. Independently installable via \texttt{pip install qsim-sdk}.
    \item \textbf{Persistence Layer}: PostgreSQL for user accounts, projects, and conversion history; Redis for caching and session management.
    \item \textbf{Observability Layer}: Health checks, Prometheus metrics, execution time histograms, and structured logging.
\end{itemize}

\subsection{Technology Stack}

\subsubsection{Frontend Technologies}
\begin{itemize}
    \item Next.js 14 with App Router for server-side rendering and routing.
    \item React 18 with TypeScript for type-safe UI development.
    \item Tailwind CSS for styling with a custom quantum-themed design.
    \item Monaco Editor for code editing with custom syntax highlighting.
    \item D3.js for 2D circuit visualization; \texttt{@react-three/fiber} for 3D rendering.
    \item Zustand for lightweight client-side state management.
    \item Chart.js for measurement histogram visualization.
\end{itemize}

\subsubsection{Backend Technologies}
\begin{itemize}
    \item Python 3.11+ for backend logic and quantum libraries.
    \item FastAPI for high-performance async API framework.
    \item Uvicorn ASGI server for production deployment.
    \item Pydantic for data validation and serialization.
    \item SQLAlchemy ORM with PostgreSQL.
    \item Redis for caching and session management.
    \item Alembic for database migrations.
\end{itemize}

\subsubsection{Quantum Computing Stack}
\begin{itemize}
    \item Cirq 1.0+ for Google quantum framework support \cite{cirq_developers_2023}.
    \item Qiskit 0.44+ for IBM quantum framework support \cite{qiskit}.
    \item PennyLane 0.33+ for Xanadu quantum framework support \cite{bergholm2018pennylane}.
    \item Custom OpenQASM 3.0 generator (\texttt{QASM3Builder}) \cite{cross2021openqasm}.
    \item QSim SDK: Custom simulation backends (statevector, density matrix, stabilizer) \cite{nielsen2010quantum}.
    \item Both SDKs deployed to PyPI for community access.
\end{itemize}

\subsection{Hybrid CPU--QPU Execution Model}

\subsubsection{Orchestration Boundary}
\begin{itemize}
    \item \textbf{CPU Responsibilities}: Circuit construction, parameter preparation, AST parsing, OpenQASM generation, validation, pre/post-processing, and result aggregation.
    \item \textbf{QSim/Simulator Responsibilities}: Gate-level circuit execution, statevector/density matrix evolution, measurement sampling, and state snapshot generation.
\end{itemize}

\subsubsection{Data Flow}
\begin{enumerate}
    \item User writes circuit in source framework (Qiskit/Cirq/PennyLane) in the Monaco editor.
    \item Frontend sends code to FastAPI backend via REST API.
    \item Backend parses code using AST or runtime execution.
    \item QCanvas SDK generates OpenQASM 3.0 intermediate representation.
    \item Validation layer checks conformance to supported subset.
    \item QSim SDK executes circuit on selected backend (statevector/density/stabilizer).
    \item Results aggregated and returned via REST/WebSocket.
    \item Frontend visualizes results (histograms, circuit diagrams, statistics).
\end{enumerate}

\section{Design Models}

\subsection{Activity Diagrams}

\subsubsection{Activity Diagram: Circuit Conversion and Execution}
This diagram presents the end-to-end flow from code input in the Monaco editor, through FastAPI conversion and validation services, into QSim execution and results visualization.
\begin{figure}[H]
    \centering
    \includegraphics[width=0.65\linewidth]{FYP-Activity-2.png}
    \caption{Activity Diagram --- Circuit Conversion and Execution}
    \label{fig:activity_convert_run}
\end{figure}

\subsubsection{Activity Diagram: User Registration}
This diagram captures the user registration and onboarding flow: a new user submits credentials, FastAPI validates and hashes the password, a PostgreSQL record is created, and a JWT session token is issued.
\begin{figure}[H]
    \centering
    \includegraphics[width=0.55\linewidth]{FYP-ActivityRegistration.png}
    \caption{Activity Diagram --- User Registration}
\end{figure}

\subsection{Sequence Diagrams}

\subsubsection{Sequence Diagram: Unified Execution Flow}
This diagram shows the interaction flow across concrete system components: the user browser, Next.js app, FastAPI routes, ConversionService/SimulationService layer, and the QSim backends.
\begin{figure}[H]
    \centering
    \includegraphics[width=\linewidth]{FYP-SequenceUnified.png}
    \caption{Sequence Diagram --- Unified Flow}
    \label{fig:seq_unified}
\end{figure}

\subsubsection{Sequence Diagram: Authentication}
\begin{figure}[H]
    \centering
    \includegraphics[width=0.95\linewidth]{FYP-SequenceAuth.png}
    \caption{Sequence Diagram --- Authentication}
\end{figure}

\subsubsection{System Sequence Diagrams (SSD)}

\textbf{SSD: Authentication}
\begin{figure}[H]
    \centering
    \includegraphics[width=0.95\linewidth]{FYP-Ssd-Auth.png}
    \caption{SSD --- Authentication}
\end{figure}

\textbf{SSD: Conversion}
\begin{figure}[H]
    \centering
    \includegraphics[width=0.92\linewidth]{FYP-Ssd-Conversion.png}
    \caption{SSD --- Conversion}
\end{figure}

\textbf{SSD: Simulation}
\begin{figure}[H]
    \centering
    \includegraphics[width=0.95\linewidth]{FYP-Ssd-Simulation.png}
    \caption{SSD --- Simulation}
\end{figure}

\clearpage
\subsection{Architecture Diagram}
This diagram provides a high-level view of the deployed QCanvas stack: the Next.js IDE, FastAPI API layer, \texttt{quantum\_converters} and \texttt{quantum\_simulator} (QSim) packages, PostgreSQL/Redis persistence, and WebSocket/REST communication channels.
\begin{figure}[H]
    \centering
    \includegraphics[width=0.95\linewidth]{FYP-SystemArchitecture.png}
    \caption{System Architecture Diagram}
    \label{fig:architecture}
\end{figure}

\clearpage
\subsection{Component Diagram}
This diagram depicts concrete module boundaries: editor and results components in the frontend, API routes and services in the backend, conversion/validation modules, QSim simulation backends, and shared utilities.
\begin{figure}[H]
    \centering
    \includegraphics[width=0.95\linewidth]{FYP-Component.png}
    \caption{Component Diagram}
    \label{fig:component}
\end{figure}

\clearpage
\subsection{Package Diagram}
This diagram outlines the actual package organization across the frontend (Next.js \texttt{app}, \texttt{components}, \texttt{lib}), backend (\texttt{app/api}, \texttt{core}, \texttt{services}, \texttt{models}), \texttt{quantum\_converters} (parsers, converters, optimizers), and \texttt{quantum\_simulator} (QSim) libraries.
\begin{figure}[H]
    \centering
    \includegraphics[width=0.95\linewidth]{Package Diagram.jpeg}
    \caption{Package Diagram}
    \label{fig:package}
\end{figure}

\clearpage
\subsection{Deployment Diagram}
This diagram shows the deployment view for production: an Nginx HTTP reverse proxy, the FastAPI backend service, the Next.js frontend, PostgreSQL/Redis containers, and containerization via Docker Compose.
\begin{figure}[H]
    \centering
    \includegraphics[width=0.95\linewidth]{FYP-Deployment.png}
    \caption{Deployment Diagram}
    \label{fig:deployment}
\end{figure}

\clearpage
\subsection{State Diagrams}

\subsubsection{State Diagram: Conversion Process}
\begin{figure}[H]
    \centering
    \includegraphics[width=0.8\linewidth]{FYP-StateConversion.png}
    \caption{State Diagram --- Conversion}
\end{figure}

\subsubsection{State Diagram: Simulation Job}
\begin{figure}[H]
    \centering
    \includegraphics[width=0.8\linewidth]{FYP-StateJob.png}
    \caption{State Diagram --- Simulation Job}
\end{figure}

\clearpage
\subsection{Data Flow Diagrams (DFD)}

\subsubsection{DFD Level 0 (Context Diagram)}
At the context level, the DFD shows QCanvas as a single system interacting with external entities: the User, Authentication Provider, and QSim backend.
\begin{figure}[H]
    \centering
    \includegraphics[width=0.8\linewidth]{FYP-DFD-lvl0.png}
    \caption{DFD Level 0}
\end{figure}

\subsubsection{DFD Level 1}
\begin{figure}[H]
    \centering
    \includegraphics[width=0.98\linewidth]{FYP-DFD-lvl1.png}
    \caption{DFD Level 1}
\end{figure}

\section{Data Design}

\subsection{Core Entities}
\begin{itemize}
    \item \textbf{Circuit}: Framework source code, parsed AST, OpenQASM 3.0 representation, metadata, and statistics.
    \item \textbf{ConversionRequest}: Source framework, target framework or QASM, source code, optimization options, validation flags.
    \item \textbf{ConversionResult}: Generated OpenQASM 3.0 code, conversion statistics, warnings, errors, execution time.
    \item \textbf{SimulationJob}: OpenQASM code, selected backend, shots, seed, execution status, progress percentage.
    \item \textbf{SimulationResult}: Measurement counts, statevector/density matrix (optional), execution time, metadata.
    \item \textbf{Project}: Files collection, active file reference, project settings, theme preferences.
    \item \textbf{User}: Authentication credentials, profile data, XP, level, achievement set.
    \item \textbf{Achievement}: Criteria type, threshold value, XP reward, unlock state.
    \item \textbf{Recipe}: Circuit snippet, framework tag, difficulty label, community tags, upvote count.
\end{itemize}

\subsection{Entity Relationship Diagram}
\begin{figure}[H]
    \centering
    \includegraphics[width=\linewidth]{FYP-ErdDiagram.png}
    \caption{Entity Relationship Diagram}
    \label{fig:erd}
\end{figure}

\subsection{Data Organization and Processing}

\begin{longtable}{p{0.22\textwidth} p{0.70\textwidth}}
\toprule
\textbf{Layer} & \textbf{Data Handled / Description} \\
\midrule
\endfirsthead
\toprule
\textbf{Layer} & \textbf{Data Handled / Description} \\
\midrule
\endhead
Input Layer & Source code (Python/QASM), examples, editor configuration (JSON), user preferences \\
Parsing Layer & Python AST, circuit operations, gate definitions, qubit mappings \\
Conversion Layer & CircuitAST intermediate format, gate mappings, OpenQASM 3.0 text, diagnostics \\
Validation Layer & Syntax errors, semantic errors, unsupported feature warnings, subset conformance checks \\
Simulation Layer & Backend selection, shots configuration, execution parameters, caching keys \\
Results Layer & Measurement counts, histograms, statevector arrays, execution metadata \\
Rendering Layer & Circuit diagrams (SVG/D3), charts (Chart.js), formatted statistics \\
Persistence Layer & Project files, conversion history, results cache, user sessions, achievements, recipes \\
\bottomrule
\caption{Data Organization and Processing}
\label{tab:data-organization}
\end{longtable}

\subsection{API Data Models}

All API endpoints use Pydantic models for validation:

\subsubsection{Request Models}
\begin{itemize}
    \item \texttt{ConversionRequest}: source\_framework, target\_framework, source\_code, style, validate
    \item \texttt{SimulationRequest}: qasm\_code, backend, shots, seed, noise\_model
    \item \texttt{FileRequest}: name, content, language, project\_id
    \item \texttt{RecipeRequest}: title, code, framework, difficulty, tags
\end{itemize}

\subsubsection{Response Models}
\begin{itemize}
    \item \texttt{ConversionResponse}: success, qasm\_code, stats, warnings, errors, timestamp
    \item \texttt{SimulationResponse}: success, results, execution\_time, backend\_info, timestamp
    \item \texttt{ErrorResponse}: error\_code, message, details, suggestions
    \item \texttt{AchievementEvent}: achievement\_id, name, xp\_awarded, new\_level
\end{itemize}
